\section{Evaluation}
\label{sec:experiments}

All formal performance results are defined on the corrected CID model retrained from the pinned iLLaDA base checkpoint~\citep{nie2026illada}. We organize evaluation around three properties that follow directly from the architecture: end-to-end tool-augmented task quality, the ability to overlap useful model computation with external latency, and the ability to assimilate delayed or changing evidence by revising already formed state. We additionally use controlled runtime interventions that require no separately retrained ablation models. Full data, optimization, checkpoint, and timing details are given in \cref{app:reproducibility}.

\subsection{Model and Evaluation Setup}
\label{sec:evaluation-setup}

The final model is trained in two stages. Stage A freezes the iLLaDA backbone and learns CID-specific channel projections, external-perception fusion, and prediction heads. Stage B then unfreezes the backbone and continues optimization under the same corrected train--runtime contract. The final checkpoint used for evaluation is fixed before benchmark results are collected; engineering smoke runs and pre-fix checkpoints are excluded from the reported comparisons.

\paragraph{Early-stage diagnostic.}
During development, an initial frozen-backbone prototype learned several downstream interaction components but failed to materialize new cognitive cells reliably in the unmodified runtime. Teacher-forced probes selected the correct source on 165/165 supervised decisions and produced strong argument representations, while positive allocation probabilities averaged only 0.110 and none of 219 positives reached the then-used 0.8 runtime threshold. A diagnostic allocation intervention that changed no model weights restored the complete need$\rightarrow$binding$\rightarrow$tool$\rightarrow$observation$\rightarrow$projection path on 29/32 held-out NQ-Open cases at the original 128-slot geometry and 30/32 with training-aligned eight-slot geometry. The initial diagnosis exposed physical-slot, logical-position, and display-shape mismatches. A subsequent contract audit additionally found inconsistent thought/display noise semantics, a self-rollout schedule that did not match runtime denoising, lifecycle/reclamation differences, unsupervised EOS-tail behavior, incompatible external-memory serialization, and an information-need ABI that could not faithfully represent concurrent bindings or explicit affected-region routing. The corrected implementation aligns these train--runtime contracts, including causal closed-loop supervision and the canonical unresolved Display boundary, uses stable multi-need slots and learned routing to affected TCT/display regions, versions the neural ABI so incompatible checkpoints fail fast, and restarts formal training from the base model. The diagnostic prototype is excluded from performance evaluation; its traces and checkpoint are retained only as development artifacts.

\subsection{End-to-End Tool-Augmented Performance}
\label{sec:end-to-end-evaluation}

The primary evaluation asks whether the corrected model can complete tool-augmented tasks reliably without diagnostic intervention. We report final-answer correctness together with tool-use precision and recall, observation coverage, stale-observation rate, convergence, and unsupported-tool-call frequency. The benchmark suite combines natural open-domain and multi-document tasks with CID-specific cases that require delayed evidence, explicit tool restraint, compositional reasoning, and persistent source use. Public-source evaluation splits are disjoint from the corresponding training splits.

Comparisons include conventional explicit-call dLLM execution and strong autoregressive tool-using systems using the same available sources where possible. For end-to-end quality comparisons, each system receives the same task input and source contents; latency-dependent experiments are separated below so hardware speed does not contaminate task correctness.

\subsection{Asynchronous Interaction Efficiency}
\label{sec:async-efficiency}

CID is intended to convert part of external waiting time into useful refinement. We therefore replay the same tasks under controlled source delays and compare end-to-end latency at matched answer quality. The main measurements are wall-clock completion time, tool-wait overlap, the amount of model work performed while I/O is outstanding, duplicate external execution, and the fraction of pre-arrival computation that remains useful after evidence arrives.

The key controlled comparison uses the same CID checkpoint in asynchronous and forced-blocking runtime modes, avoiding a retraining confound. We additionally compare against blocking and asynchronous autoregressive tool execution when an equivalent source interface is available. Quality comparisons use a logical event clock, while wall-clock results are reported on matched hardware.

\subsection{Continuous Assimilation and Revision}
\label{sec:continuous-revision}

To test the specifically diffusion-native claim, we construct cases in which useful evidence arrives after partial cognition or display content has formed, changes during generation, or is delivered incrementally. Static-source cases distinguish one-time injection from persistent cognitive re-projection. Dynamic-source cases change the authoritative value during the trajectory, and streaming cases reveal increasingly informative evidence chunks.

We measure stale-value rate, observation-to-correction lag, affected-region routing precision/recall, the fraction of routed thought/display regions revised after arrival, preservation of unrelated state, and final correctness. A successful system should identify and revise evidence-dependent regions without restarting the trajectory or repeatedly executing an unchanged source.

\subsection{Ablation Study}
\label{sec:ablations}

Because CID's learned channels, interaction heads, and runtime state transitions are tightly coupled, removing most learned mechanisms would require training a different model and would confound the effect of the mechanism with a new optimization run. We therefore use three targeted ablations that cover the main interaction claims while keeping the study compact. One controlled development ablation compares two models trained from the same base model, data materialization, contract, batch geometry, and optimization budget, differing only in whether interaction needs are decoded from the provisional updated thought state. The other two interventions reuse the final corrected checkpoint and alter only runtime execution policy.

\begin{table}[t]
  \centering
  \caption{Three targeted CID ablations. The provisional-state comparison uses matched epoch-2 checkpoints at 62,645 optimizer steps and fixed-seed free-rollout validation. The two runtime ablations reuse the final corrected checkpoint.}
  \label{tab:cid-ablation}
  \scriptsize
  \setlength{\tabcolsep}{2.5pt}
  \begin{tabular}{p{0.19\linewidth}p{0.22\linewidth}p{0.28\linewidth}p{0.22\linewidth}}
    \toprule
    Ablation & Comparison & Primary metrics & Result \\
    \midrule
    Provisional-state decoding & stale versus provisional thought & display exact; Need F1; recovery failure & 13.6$\rightarrow$22.8; 1.9$\rightarrow$9.1; 8.14$\rightarrow$5.81 \\
    Asynchronous execution & async versus forced blocking & latency; wait overlap; task quality & \textit{pending runtime replay} \\
    Continuous revision & normal versus freeze after materialization & final accuracy; stale rate; repair rate & \textit{pending runtime replay} \\
    \bottomrule
  \end{tabular}
\end{table}

\paragraph{Provisional-state-aware interaction decoding.}
The first ablation comes from a preserved controlled retraining pair rather than from the final corrected model. Both runs use the same frozen CID-Dataset materialization and neural contract, 422,230 training examples, 3,146,786 transitions, 44.29M trainable parameters, effective batch size 96, and the same validation seed. The only mechanism change makes need, argument, refresh, and routing predictions condition on the provisional thought state after the current thought update, rather than on the stale pre-update hidden state. At the common epoch-2 budget, this change raises free-rollout materialized-display exact rate from 13.6\% to 22.8\% and display-token accuracy from 20.4\% to 27.7\%. Need F1 rises from 1.9\% to 9.1\%, while rollout recovery failure falls from 8.14\% to 5.81\%. We report this pair only as an isolated mechanism ablation: the later corrected CID model contains additional train--runtime fixes and is evaluated separately in the main results.

\paragraph{Asynchronous execution.}
The second ablation will run the final corrected checkpoint in its normal asynchronous runtime and in a forced-blocking mode that prevents useful denoising/refinement while an external request is outstanding. The comparison keeps model weights, requests, source results, and latency schedules fixed. We will report end-to-end latency, tool-wait overlap, useful pre-arrival computation, and final task quality. This directly tests whether non-blocking interaction contributes measurable efficiency rather than merely changing the execution trace.

\paragraph{Continuous revision.}
The third ablation will compare the normal runtime against a freeze-after-materialization policy using the same final checkpoint. In the ablated runtime, a cognitive/display region is frozen after its first materialization and cannot be revised by later diffusion iterations or newly arrived evidence. We will measure final correctness, stale-value rate, observation-to-correction lag, and the rate at which initially incorrect materialized content is repaired. This isolates the value of continuous revision without retraining a model whose learned representation was optimized for a different state topology.

\subsection{Controlled Runtime Analysis}
\label{sec:runtime-analysis}

Beyond the two runtime ablations above, we analyze persistent re-projection versus one-time percept injection, need-conditioned refresh versus no refresh, cache reuse versus forced re-execution, and alternative source-latency schedules. These sensitivity and runtime controls reuse the same final checkpoint and are kept separate from the three headline ablations in \cref{tab:cid-ablation}.

The broader CID hypothesis remains falsifiable. The architecture is not supported if the corrected model cannot bind tools reliably, if asynchronous computation is mostly discarded after observations arrive, if revision damages more correct state than it repairs, or if simpler explicit-call or asynchronous autoregressive systems match quality and latency with lower complexity.
